\documentclass[10pt,a4paper]{amsart}
\usepackage[margin=19mm]{geometry}
\usepackage{amsmath,amssymb,mathtools}
\usepackage[scr=rsfs,cal=euler]{mathalfa}
\usepackage{xfrac}
\usepackage[unicode,hidelinks,bookmarks=false]{hyperref}
\newcommand{\ie}{i.e.\ }
\newcommand{\cf}{cf.\ }
\newcommand{\resp}{respectively\ }
\newcommand{\Subsection}{\subsection}
\providecommand{\nobreakdash}{\nobreak\hbox{-}}
\setlength{\emergencystretch}{2em}
\multlinegap0pt
\usepackage{amssymb}
\usepackage{xcolor}
\newtheorem{lemma}{Lemma}
\title{Semi--finite Vector Bundles on Complex Tori}
\author{Pavan Adroja, Sanjay Amrutiya}
\date{Source: arXiv:2603.17661v1 (CC BY 4.0)}
\begin{document}
\maketitle

\noindent\emph{Source and licence.} Semi--finite Vector Bundles on Complex Tori. Version arXiv:2603.17661v1 (CC BY 4.0), distributed by arXiv under the Creative Commons Attribution 4.0 International licence, http://creativecommons.org/licenses/by/4.0/, as recorded in the arXiv licence metadata for this exact version and shown on the abstract page https://arxiv.org/abs/2603.17661v1. Full licence notice: \texttt{licenses/arxiv-2603.17661-CC-BY-4.0.txt}.\par\smallskip

\noindent\textit{Editorial scope.} Counted unit: the article's lemma sf=f+u (source0.tex lines 182--186). The article's complete original proof of it is delivered with it (source0.tex lines 187--239, one \texttt{proof} environment of 53 lines). No mathematics is written, repaired or re-derived: the statement and the proof are the article's own lines, in the article's own notation, and no retained line is substituted or re-flowed.  Retained uncounted as the source-local context the proof needs: lines 146--150; lines 166--169.  Nothing was omitted from any retained span, so the package carries no omission record.  No retained line contains a citation, so the wrapper carries no bibliography and no bibliography entry.  No retained line contains a figure environment, an image-inclusion command or drawing code, so no image is delivered and the package declares no asset dependency.  Editorial formatting only, in addition: an AMS \texttt{amsart} preamble that restates the article's own declarations and macros used by the retained lines, namely the environments \texttt{lemma} on separate counters, together with the standard packages those lines need.  The retained environments print in this document as Lemma 1 in order of appearance, whereas the article prints the same items with the numbering of its own full document.  The complete native source of the article as distributed in the arXiv e-print for this exact version is bundled unchanged as \texttt{sources/196638/source0.tex}.
\begin{lemma}\label{finite is sum of torsion on torus}
	Let $X=V/\Lambda$ be a complex torus of complex dimension $g$, and let $E$ be a finite holomorphic vector bundle on $X$. Then, 
	$$E\simeq \bigoplus_{i=1}^n L_i\,,$$
	where $L_i$'s are torsion line bundles on $X$.
\end{lemma}

\begin{lemma}\label{H^0=H^1=0}
	Let $X$ be a complex torus of complex dimension $g$, and let $L$ be a non-trivial torsion line bundle on $X$. Then, 
	$$\mathrm{H}^0(X,L)=\mathrm{H}^1(X,L)=0.$$
\end{lemma}

\begin{lemma}\label{sf=f+u}
	Let $X$ be a complex torus of complex dimension $g$, and let $E$ be a semi--finite holomorphic vector bundle on $X$. Then, 
	$$E\simeq \bigoplus_{i=1}^n U_i \otimes L_i\,, $$
	where $L_i$'s are torsion line bundles on $X$ and $U_i$'s are unipotent bundles on $X$.
\end{lemma}

\begin{proof}
	We proceed by induction on the rank of \(E\). If \(\mathrm{rank}(E)=1\), then by the definition of semi--finite bundles and Lemma~\ref{finite is sum of torsion on torus}, \(E\) must be a torsion line bundle. In this case, the statement holds with \(U=\mathcal{O}_X\) and \(L=E\). Assume the statement holds for all semi--finite bundles of rank \(< r\), and let \(\mathrm{rank}(E)=r\). By Lemma~\ref{finite is sum of torsion on torus} and the definition of semi--finite bundles, there exists a filtration
	\[
	0 = E_0 \subset E_1 \subset \cdots \subset E_{r-1} \subset E_r = E
	\]
	by subbundles such that each successive quotient \(E_i/E_{i-1}\) is a torsion line bundle. In particular, we have a short exact sequence
	\begin{equation}\label{eq:main-sequence}
		0 \longrightarrow E_{r-1} \longrightarrow E \longrightarrow L \longrightarrow 0 \,,
	\end{equation}
	where \(L = E_r/E_{r-1}\) is a torsion line bundle. Since \(\mathrm{rank}(E_{r-1})=r-1\), by the induction hypothesis, we have
	\[
	E_{r-1} \simeq \bigoplus_i U_i \otimes L_i \,,
	\]
	where each \(U_i\) is unipotent and each \(L_i\) is a torsion line bundle on \(X\). Split the direct sum for $E_{r-1}$ according to whether $L_j\simeq L$ or $L_j \not \simeq L$. Write
	$$E_{r-1} \simeq W \oplus (U_0 \otimes L) \,,$$
	where $W:= \bigoplus_{j:L_j \not \simeq L} U_j \otimes L_j$ and $U_0:= \bigoplus_{j:L_j \simeq L} U_j$. Note that $U_0$ is unipotent as it is direct sum of unipotent bundles.
	
	
	Since 
	$$\mathrm{Ext}^1 \bigl( L, E_{r-1} \bigr) \simeq \mathrm{Ext}^1 \bigl( L, W \bigr) \oplus \mathrm{Ext}^1 \bigl( L, U_0 \otimes L \bigr)\,,$$
	we have $E \simeq W' \oplus E'$ such that $W'$ and $E'$ fit into the following exact sequences: 
	\begin{equation}\label{E:W'}
		0\longrightarrow W \longrightarrow W' \longrightarrow L \longrightarrow 0 \,.
	\end{equation}
	\begin{equation}\label{E:E'}
		0\longrightarrow U_0 \otimes L \longrightarrow E' \longrightarrow L \longrightarrow 0 \,.
	\end{equation}
	Tensoring the exact sequence \eqref{E:E'} with \(L^{-1}\), and using that tensoring by a line bundle is an exact auto-equivalence of the category of vector bundles on \(X\), we obtain an equivalent extension
	\begin{equation}\label{eq:quence}
		0 \longrightarrow U_0 \longrightarrow E'\otimes L^{-1} \longrightarrow \mathcal{O}_X \longrightarrow 0 \,.
	\end{equation}
	From the above sequence, $U':=E'\otimes L^{-1}$ is a unipotent bundle. Hence, $E'=U' \otimes L$. Therefore, it suffices to show that the exact sequence \eqref{E:W'} splits. The sequence \eqref{E:W'} splits if the sequence 
	$$0 \longrightarrow U_j \otimes L_j \longrightarrow F_j \longrightarrow L \longrightarrow 0 $$
	splits for every $j$, where $L_j \not \simeq L$. Thus it is enough to prove that any extension
	\begin{equation}\label{eq:Ui-extension}
		0 \longrightarrow U \otimes L' \longrightarrow F \longrightarrow \mathcal{O}_X \longrightarrow 0
	\end{equation}
	splits, where \(U\) is a unipotent bundle and \(L'\) is a non-trivial torsion line bundle on \(X\).
	
	Since \(U\) is unipotent, it admits a filtration
	\[
	0 = U_0 \subset U_1 \subset \cdots \subset U_{m-1} \subset U_m = U
	\]
	with successive quotients \(U_j/U_{j-1} \simeq \mathcal{O}_X\). Tensoring by \(L'\) gives
	\[
	0 \longrightarrow U_{m-1}\otimes L' \longrightarrow U \otimes L' \longrightarrow L' \longrightarrow 0 \,.
	\]
	By induction on \(\mathrm{rank}(U)\), it suffices to show that any extension
	\begin{equation}\label{eq:L-extension}
		0 \longrightarrow L' \longrightarrow G \longrightarrow \mathcal{O}_X \longrightarrow 0
	\end{equation}
	splits. Since $\mathrm{Ext}^1(\mathcal{O}_X, L') \simeq \mathrm{H}^1(X, L')$ and $\mathrm{H}^1(X, L')=0$ by Lemma \ref{H^0=H^1=0}, the sequence (\ref{eq:L-extension}) splits.
\end{proof}

\end{document}
